\newpage
\phantomsection
\label{prf:differentiability-and-one-sided-derivatives}

\begin{remark*}[Return]
\hyperref[thm:differentiability-and-one-sided-derivatives]{Return to Theorem}
\end{remark*}

\begin{theorem*}[Differentiability and One-Sided Derivatives]
Let $I \subseteq \mathbb{R}$ be an open interval and let $f : I \to \mathbb{R}$. Fix $c \in I$. Then $f$ is differentiable at $c$ if and only if both one-sided derivatives $f'_-(c)$ and $f'_+(c)$ exist and are equal, in which case $f'(c) = f'_-(c) = f'_+(c)$.
\end{theorem*}

\begin{proof}
\textbf{Professional Standard Proof.}~\\
($\Rightarrow$) If $f$ is differentiable at $c$, then by definition the two-sided limit $\lim_{x \to c} \frac{f(x) - f(c)}{x - c} = f'(c)$ exists. By the characterization of two-sided limits in terms of one-sided limits, both $f'_-(c)$ and $f'_+(c)$ exist and equal $f'(c)$.

($\Leftarrow$) If both $f'_-(c)$ and $f'_+(c)$ exist and equal some $L \in \mathbb{R}$, then both one-sided limits of the difference quotient exist and equal $L$. By the characterization of two-sided limits, the two-sided limit $\lim_{x \to c} \frac{f(x) - f(c)}{x - c} = L$ exists, so $f$ is differentiable at $c$ with $f'(c) = L$.
\end{proof}

\begin{proof}
\textbf{Detailed Learning Proof.}~\\

\textbf{Step 1.} Establish the forward implication: differentiability implies equal one-sided derivatives.

Assume $f$ is differentiable at $c$. By the definition of differentiability, the two-sided limit
\[
\lim_{x \to c} \frac{f(x) - f(c)}{x - c} = f'(c)
\]
exists. The fundamental theorem relating two-sided and one-sided limits states that a two-sided limit exists if and only if both corresponding one-sided limits exist and are equal. Applying this characterization to the difference quotient, both
\[
f'_-(c) = \lim_{x \to c^-} \frac{f(x) - f(c)}{x - c} \quad \text{and} \quad f'_+(c) = \lim_{x \to c^+} \frac{f(x) - f(c)}{x - c}
\]
exist and satisfy $f'_-(c) = f'_+(c) = f'(c)$.

\textbf{Step 2.} Establish the reverse implication: equal one-sided derivatives imply differentiability.

Assume both one-sided derivatives $f'_-(c)$ and $f'_+(c)$ exist and equal some common value $L \in \mathbb{R}$. This means
\[
\lim_{x \to c^-} \frac{f(x) - f(c)}{x - c} = L \quad \text{and} \quad \lim_{x \to c^+} \frac{f(x) - f(c)}{x - c} = L.
\]
Since both one-sided limits of the difference quotient exist and have the same value $L$, the characterization of two-sided limits guarantees that the two-sided limit exists and equals $L$:
\[
\lim_{x \to c} \frac{f(x) - f(c)}{x - c} = L.
\]
By the definition of differentiability, this establishes that $f$ is differentiable at $c$ with derivative $f'(c) = L = f'_-(c) = f'_+(c)$.
\end{proof}

\begin{remark*}[Proof structure]
This is a pure equivalence proof that reduces entirely to the characterization of two-sided limits in terms of one-sided limits. The derivative is defined as a two-sided limit of the difference quotient, so the equivalence between differentiability and equal one-sided derivatives follows immediately from applying the one-sided limit criterion to this defining limit. No calculation or construction is required — the result is a direct translation of the limit-theoretic characterization.
\end{remark*}

\begin{remark*}[Dependencies]~\\
\begin{itemize}
  \item \hyperref[def:derivative-at-point]{Definition of Differentiability at a Point}
  \item \hyperref[thm:two-sided-one-sided-limits]{Two-Sided Limit and One-Sided Limits}
\end{itemize}
\end{remark*}